\chapter{2009年新课标卷（理）}

\section{单选题}

\begin{problem}
已知集合 \(A = \{1, 3, 5, 7, 9\}\)，\(B = \{0, 3, 6, 9, 12\}\)，则 \(A \cap \complement_{\N} B =\)（\quad）

\choices
    {\(\{1, 5, 7\}\)}
    {\(\{3, 5, 7\}\)}
    {\(\{1, 3, 9\}\)}
    {\(\{1, 2, 3\}\)}
\end{problem}

\begin{problem}
复数 \(\frac{3 + 2\i}{2 - 3\i} - \frac{3 - 2\i}{2 + 3\i} =\)（\quad）

\choices
    {\(0\)}
    {\(2\)}
    {\(-2\i\)}
    {\(2\)}
\end{problem}

\begin{problem}
对变量 \(x\)，\(y\) 有观测数据 \((x_{i}, y_{i})\)（\(i = 1, 2, \dots, 10\)），得散点图 1；对变量 \(u\)，\(v\) 有观测数据 \((u_{i}, v_{i})\)（\(i = 1, 2, \dots, 10\)），得散点图 2. 由这两个散点图可以判断（\quad）

\bitmapfigure[max width=0.82\linewidth]{img/2009/new_curriculum_science/q3_fig1.png}

\choices
    {变量 \(x\) 与 \(y\) 正相关，\(u\) 与 \(v\) 正相关}
    {变量 \(x\) 与 \(y\) 正相关，\(u\) 与 \(v\) 负相关}
    {变量 \(x\) 与 \(y\) 负相关，\(u\) 与 \(v\) 正相关}
    {变量 \(x\) 与 \(y\) 负相关，\(u\) 与 \(v\) 负相关}
\end{problem}

\begin{problem}
双曲线 \(\frac{x^2}{4} - \frac{y^2}{12} = 1\) 的焦点到渐近线的距离为（\quad）

\choices
    {\(2 \sqrt{3}\)}
    {\(2\)}
    {\(\sqrt{3}\)}
    {\(1\)}
\end{problem}

\begin{problem}
有四个关于三角函数的命题：\(p_1: \exists x \in \R\)，\(\sin^2 \frac{x}{2} + \cos^2 \frac{x}{2} = \frac{1}{2};\) \(p_2: \exists x\)，\(y \in \R\)，\(\sin (x - y) = \sin x - \sin y.\) \(p_3\colon \forall x \in \left[0,\pi \right]\)，\(\sqrt{\frac{1 - \cos 2x}{2}} = \sin x;\) \(p_4\colon \sin x = \cos y \Rightarrow x + y = \frac{\pi}{2}\). 其中假命题的是（\quad）

\choices
    {\(p_1, p_4\)}
    {\(p_2\)，\(p_4\)}
    {\(p_1\)，\(p_3\)}
    {\(p_2\)，\(p_3\)}
\end{problem}

\begin{problem}
设 \(x\)，\(y\) 满足 \(\begin{cases}
2x + y \geqslant 4,\\
x - y \geqslant -1,\\
x - 2y \leqslant 2,
\end{cases}\) 则 \(z = x + y\)（\quad）

\choices
    {有最小值 \(2\)，最大值 \(3\)}
    {有最小值 \(2\)，无最大值}
    {有最大值 \(3\)，无最小值}
    {既无最小值，也无最大值}
\end{problem}

\begin{problem}
等比数列 \(\{a_{n}\}\) 的前 \(n\) 项和为 \(S_{n}\)，且 \(4a_{1}\)，\(2a_{2}\)，\(a_{3}\) 成等差数列. 若 \(a_{1} = 1\)，则 \(S_{4} =\)（\quad）

\choices
    {\(7\)}
    {\(8\)}
    {\(15\)}
    {\(16\)}
\end{problem}

\begin{problem}
如图，正方体 \(ABCD - A_{1}B_{1}C_{1}D_{1}\) 的棱长为 \(1\)，线段 \(B_{1}D_{1}\) 上有两个动点 \(E\)，\(F\)，且 \(EF = \frac{\sqrt{2}}{2}\)，则下列结论中错误的是（\quad）

\bitmapfigure[max width=0.42\linewidth]{img/2009/new_curriculum_science/q8_fig1.png}

\choices
    {\(AC \perp BE\)}
    {\(EF //\) 平面 \(ABCD\)}
    {三棱锥 \(A - BEF\) 的体积为定值}
    {异面直线 \(AE\)，\(BF\) 所成的角为定值}
\end{problem}

\begin{problem}
已知 \(O\)，\(N\)，\(P\) 在 \(\triangle ABC\) 所在平面内，且 \(\left|\overrightarrow{OA}\right| = \left|\overrightarrow{OB}\right| = \left|\overrightarrow{OC}\right|\)，\(\overrightarrow{NA} + \overrightarrow{NB} + \overrightarrow{NC} = \mathbf{0}\)，\(\overrightarrow{PA} \cdot \overrightarrow{PB} = \overrightarrow{PB} \cdot \overrightarrow{PC} = \overrightarrow{PC} \cdot \overrightarrow{PA}\)，则点 \(O\)，\(N\)，\(P\) 依次是 \(\triangle ABC\) 的（\quad）

\choices
    {重心，外心，垂心}
    {重心，外心，内心}
    {外心，重心，垂心}
    {外心，垂心，内心}
\end{problem}

\begin{problem}
如果执行如图的程序框图，输入 \(x = -2\)，\(h = 0.5\)，那么输出的各个数的和等于（\quad）

\texfigure[max width=0.58\linewidth]{img_repaint/2009/new_curriculum_science/q10_fig1.tex}

\choices
    {\(3\)}
    {\(3.5\)}
    {\(4\)}
    {\(4.5\)}
\end{problem}

\begin{problem}
一个棱锥的三视图如图，则该棱锥的全面积（单位：\(\mathrm{cm}^2\)）为（\quad）

\bitmapfigure[max width=0.48\linewidth]{img/2009/new_curriculum_science/q11_fig1.png}

\choices
    {\(48 + 12\sqrt{2}\)}
    {\(48 + 24\sqrt{2}\)}
    {\(36 + 12\sqrt{2}\)}
    {\(36 + 24\sqrt{2}\)}
\end{problem}

\begin{problem}
用 \(\min \{a, b, c\}\) 表示 \(a\)，\(b\)，\(c\) 三个数中的最小值. 设 \(f(x) = \min \{2^x, x + 2, 10 - x\}\)（\(x \geqslant 0\)），则 \(f(x)\) 的最大值为（\quad）

\choices
    {\(4\)}
    {\(5\)}
    {\(6\)}
    {\(7\)}
\end{problem}

\section{填空题}

\begin{problem}
设已知抛物线 \(C\) 的顶点在坐标原点，焦点为 \(F(1,0)\)，直线 \(l\) 与抛物线 \(C\) 相交于 \(A\)，\(B\) 两点. 若 \(AB\) 的中点为 \((2,2)\)，则直线 \(l\) 的方程为\fillinblank{}.
\end{problem}

\begin{problem}
已知函数 \( y = \sin (\omega x + \varphi) \) （\(\omega > 0, -\pi \leqslant \varphi < \pi\)）的图象如图所示，则 \(\varphi = \)\fillinblank{}.

\bitmapfigure[max width=0.40\linewidth]{img/2009/new_curriculum_science/q14_fig1.png}
\end{problem}

\begin{problem}
\(7\)名志愿者中安排\(6\)人在周六，周日两天参加社区公益活动. 若每天安排\(3\)人，则不同的安排方案共有\fillinblank{} 种.（用数字作答）
\end{problem}

\begin{problem}
等差数列 \(\{a_{n}\}\) 前 \(n\) 项和为 \(S_{n}\). 已知 \(a_{m-1} + a_{m+1} - a_{m}^{2} = 0\)，\(S_{2m-1} = 38\)，则 \(m =\)\fillinblank{}.
\end{problem}

\section{解答题}

\begin{problem}
为了测量两山顶 \(M\)，\(N\) 间的距离，飞机沿水平方向在 \(A\)，\(B\) 两点进行测量. \(A\)，\(B\)，\(M\)，\(N\) 在同一个铅垂平面内（如示意图），飞机能够测量的数据有俯角和 \(A\)，\(B\) 间的距离，请设计一个方案，包括：

\begin{enumerate}
\item 指出需要测量的数据（用字母表示，并在图中标出）；
\item 用文字和公式写出计算 \(M\)，\(N\) 间距离的步骤.
\end{enumerate}

\bitmapfigure[max width=0.42\linewidth]{img/2009/new_curriculum_science/q17_fig1.png}
\end{problem}

\begin{problem}
某工厂有工人 \(1000\) 名，其中 \(250\) 名工人参加过短期培训（称为 \(A\) 类工人），另外 \(750\) 名工人参加过长期培训（称为 \(B\) 类工人）.现用分层抽样方法（按 \(A\) 类，\(B\) 类分二层）从该工厂的工人中共抽查 \(100\) 名工人，调查他们的生产能力（此处生产能力指一天加工的零件数）.

\begin{enumerate}
\item 求甲，乙两工人都被抽到的概率，其中甲为 \(A\) 类工人，乙为 \(B\) 类工人；
\item 从 \(A\) 类工人中的抽查结果和从 \(B\) 类工人中的抽查结果分别如下表 1 和表 2.

\begin{center}
表 1

\begin{tabular}{c|ccccc}
生产能力分组 & \([100,110)\) & \([110,120)\) & \([120,130)\) & \([130,140)\) & \([140,150)\) \\
\hline
人数 & \(4\) & \(8\) & \(x\) & \(5\) & \(3\)
\end{tabular}
\end{center}

\begin{center}
表 2

\begin{tabular}{c|cccc}
生产能力分组 & \([110,120)\) & \([120,130)\) & \([130,140)\) & \([140,150)\) \\
\hline
人数 & \(6\) & \(y\) & \(36\) & \(18\)
\end{tabular}
\end{center}

\begin{enumerate}
\item 先确定 \(x\)，\(y\)，再完成下列频率分布直方图. 就生产能力而言，\(A\) 类工人中个体间的差异程度与 \(B\) 类工人中个体间的差异程度哪个更小？（不用计算，可通过观察直方图直接回答结论）
\begin{center}
\texfigure[max width=0.90\linewidth]{img_repaint/2009/new_curriculum_science/q18_fig1.tex}
\end{center}
\item 分别估计 \(A\) 类工人和 \(B\) 类工人生产能力的平均数，并估计该工厂工人的生产能力的平均数.（同一组中的数据用该组区间的中点值作代表）
\end{enumerate}
\end{enumerate}
\end{problem}

\begin{problem}
如图，四棱锥 \(S - ABCD\) 的底面是正方形，每条侧棱的长都是地面边长的 \(\sqrt{2}\) 倍，\(P\) 为侧棱 \(SD\) 上的点.

\begin{enumerate}
\item 求证：\(AC \perp SD\)；
\item 若 \(SD \perp\) 平面 \(PAC\)，求二面角 \(P - AC - D\) 的大小；
\item 在上一问的条件下，侧棱 \(SC\) 上是否存在一点 \(E\)，使得 \(BE //\) 平面 \(PAC\). 若存在，求 \(SE:EC\) 的值；若不存在，试说明理由.
\end{enumerate}

\bitmapfigure[max width=0.34\linewidth]{img/2009/new_curriculum_science/q19_fig1.png}

\end{problem}

\begin{problem}
已知椭圆 \(C\) 的中心为直角坐标系 \(xOy\) 的原点，焦点在 \(x\) 轴上，它的一个顶点到两个焦点的距离分别是 \(7\) 和 \(1\).

\begin{enumerate}
\item 求椭圆 \(C\) 的方程；
\item 若 \(P\) 为椭圆 \(C\) 上的动点，\(M\) 为过 \(P\) 且垂直于 \(x\) 轴的直线上的点，\(\frac{|OP|}{|OM|} = \lambda\)（\(\lambda\) 为椭圆 \(C\) 的离心率）. 求点 \(M\) 的轨迹方程，并说明轨迹是什么曲线.
\end{enumerate}
\end{problem}

\begin{problem}
已知函数 \( f(x) = (x^3 + 3x^2 + ax + b)\e^{-x} \).

\begin{enumerate}
\item 若 \(a = b = -3\)，求 \(f(x)\) 的单调区间；
\item 若 \(f(x)\) 在 \((-\infty, \alpha)\)，\((2,\beta)\) 单调增加，在 \((\alpha, 2)\)，\((\beta, +\infty)\) 单调减少，证明：\(\beta - \alpha < 6\).
\end{enumerate}
\end{problem}

\begin{problem}
如图，已知 \(\triangle ABC\) 的两条角平分线 \(AD\) 和 \(CE\) 相交于 \(H\)，\(\angle B = 60^{\circ}\)，\(F\) 在 \(AC\) 上，且 \(AE = AF\).

\begin{enumerate}
\item 证明：\(B\)，\(D\)，\(H\)，\(E\) 四点共圆；
\item 证明：\(CE\) 平分 \(\angle DEF\).
\end{enumerate}

\bitmapfigure[max width=0.34\linewidth]{img/2009/new_curriculum_science/q22_fig1.png}

\end{problem}

\begin{problem}
已知曲线 \(C_{1}: \begin{cases}
x = -4 + \cos t,\\
y = 3 + \sin t,
\end{cases}\)（\(t\) 为参数），\(C_{2}: \begin{cases}
x = 8\cos \theta,\\
y = 3\sin \theta.
\end{cases}\)（\(\theta\) 为参数）.

\begin{enumerate}
\item 化 \(C_1\)，\(C_2\) 的方程为普通方程，并说明它们分别表示什么曲线；
\item 若 \(C_1\) 上的点 \(P\) 对应的参数为 \(t = \frac{\pi}{2}\)，\(Q\) 为 \(C_2\) 上的动点，求 \(PQ\) 中点 \(M\) 到直线 \(C_{3}: \begin{cases}
x = 3 + 2t,\\
y = -2 + t,
\end{cases}\)（\(t\) 为参数）距离的最小值.
\end{enumerate}
\end{problem}

\begin{problem}
如图，\(O\) 为数轴的原点，\(A\)，\(B\)，\(M\) 为数轴上三点，\(C\) 为线段 \(OM\) 上的动点，设 \(x\) 表示 \(C\) 与原点的距离，\(y\) 表示 \(C\) 到 \(A\) 距离 \(4\) 倍与 \(C\) 到 \(B\) 距离的 \(6\) 倍的和.

\begin{enumerate}
\item 将 \(y\) 表示成 \(x\) 的函数；
\item 要使 \(y\) 的值不超过 \(70\)，\(x\) 应该在什么范围内取值.
\end{enumerate}

\bitmapfigure[max width=0.44\linewidth]{img/2009/new_curriculum_science/q24_fig1.png}

\end{problem}
